% Lecture 11 - Page 6 (13-Dec-1994).
% Fragment for % Lecture 11 - Page 6 (13-Dec-1994).
% Fragment for % Lecture 11 - Page 6 (13-Dec-1994).
% Fragment for % Lecture 11 - Page 6 (13-Dec-1994).
% Fragment for \input{courses/pages/curs11-pagina6-en}.
% No preamble or forced page break.
% Faithful reconstruction; uncertain readings are recorded in comments.

\subsection*{Lecture 11 - Page 6}

% Continuation of the isomorphism theorem proof from page 5.
\noindent $\underbrace{\bar a}_{aH}\in\ker f^*
\iff aH=aN\iff a\in H$.
% In the first line, aH=aN and a\in H are the apparent readings,
% but their relationship to the preceding homomorphism is unclear.
\[
\ker f^*=\{x\mid (\ldots)\}=\frac{H}{N}.
\]
% The expression after the vertical bar is partly cancelled/illegible.
% H/N appears to be written in the original. It has not been silently
% corrected to an alternative mathematical expression.

\bigskip
\noindent\underline{Applications:} \quad $(\mathbb Z,+,0)$.

\medskip
\noindent $H\leq\mathbb Z\ \Longrightarrow\ (\exists!)\,u\in\mathbb N^*$
 such that $H=u\mathbb Z$.
% The manuscript appears to specify N^*. Taken literally this excludes
% the trivial subgroup H={0}, for which the generator is 0.

\bigskip
\noindent Let $m,n\in\mathbb N^*$. I show that
\[
mn=[m,n]\cdot(m,n).
\]
% Brackets denote least common multiple; parentheses greatest common divisor.
\[
m\mathbb Z+n\mathbb Z\leq\mathbb Z
\ \Longrightarrow\ (\exists!)\,d\in\mathbb N\ \text{such that}\
m\mathbb Z+n\mathbb Z=d\mathbb Z.
\]
\noindent I show that $d=(m,n)$.

\medskip
\noindent There exist $x,y\in\mathbb Z$ such that
\[
mx+ny=d
\quad\Longrightarrow\quad
\text{if } c\mid m\ \text{and}\ c\mid n,
\ \text{then}\ c\mid d.
\]
\[
m\mathbb Z\subseteq d\mathbb Z,\qquad
m\in m\mathbb Z\ \Longrightarrow\ (\exists)\,q\in\mathbb Z
\ \text{such that}\ m=dq,\quad d\mid m.
\]
\noindent Similarly, $d\mid n$ $\Longrightarrow d=(m,n)$.

\bigskip
\noindent $m\mathbb Z\cap n\mathbb Z\leq\mathbb Z\ \Longrightarrow\
(\exists!)\,u\in\mathbb N$ such that
\[
m\mathbb Z\cap n\mathbb Z=u\mathbb Z.
\]
\noindent $u\in m\mathbb Z$ and $u\in n\mathbb Z$
$\Longrightarrow m\mid u$ and $n\mid u$.

\noindent If $m\mid b$ and $n\mid b$,
$\stackrel{?}{\Longrightarrow} b\mid u$.
% The handwritten final relation appears to be b|u, although the
% subsequent membership statement implies u|b. Check the original.
\[
\underbrace{b\in m\mathbb Z}_{m\mid b},\quad
\underbrace{b\in n\mathbb Z}_{n\mid b}
\quad\Longrightarrow\quad b\in u\mathbb Z.
\]

\bigskip
\noindent First Isomorphism Theorem $\Longrightarrow$
\[
\frac{m\mathbb Z+n\mathbb Z}{n\mathbb Z}
\cong \frac{m\mathbb Z}{m\mathbb Z\cap n\mathbb Z}.
\]
\[
\Longrightarrow\qquad
\underbrace{\frac{d\mathbb Z}{n\mathbb Z}}_{\cong\,\mathbb Z_{n_1}}
\cong
\underbrace{\frac{m\mathbb Z}{u\mathbb Z}}_{\cong\,\mathbb Z_{u_1}}.
\]
% The n_1 and u_1 subscripts are slightly faded in the manuscript.
\[
u=m u_1,\qquad n=d n_1.
\]
\[
u_1=n_1\quad\Longrightarrow\quad
\frac{u}{m}=\frac{n}{d}\quad\Longrightarrow\quad du=mn.
\]
.
% No preamble or forced page break.
% Faithful reconstruction; uncertain readings are recorded in comments.

\subsection*{Lecture 11 - Page 6}

% Continuation of the isomorphism theorem proof from page 5.
\noindent $\underbrace{\bar a}_{aH}\in\ker f^*
\iff aH=aN\iff a\in H$.
% In the first line, aH=aN and a\in H are the apparent readings,
% but their relationship to the preceding homomorphism is unclear.
\[
\ker f^*=\{x\mid (\ldots)\}=\frac{H}{N}.
\]
% The expression after the vertical bar is partly cancelled/illegible.
% H/N appears to be written in the original. It has not been silently
% corrected to an alternative mathematical expression.

\bigskip
\noindent\underline{Applications:} \quad $(\mathbb Z,+,0)$.

\medskip
\noindent $H\leq\mathbb Z\ \Longrightarrow\ (\exists!)\,u\in\mathbb N^*$
 such that $H=u\mathbb Z$.
% The manuscript appears to specify N^*. Taken literally this excludes
% the trivial subgroup H={0}, for which the generator is 0.

\bigskip
\noindent Let $m,n\in\mathbb N^*$. I show that
\[
mn=[m,n]\cdot(m,n).
\]
% Brackets denote least common multiple; parentheses greatest common divisor.
\[
m\mathbb Z+n\mathbb Z\leq\mathbb Z
\ \Longrightarrow\ (\exists!)\,d\in\mathbb N\ \text{such that}\
m\mathbb Z+n\mathbb Z=d\mathbb Z.
\]
\noindent I show that $d=(m,n)$.

\medskip
\noindent There exist $x,y\in\mathbb Z$ such that
\[
mx+ny=d
\quad\Longrightarrow\quad
\text{if } c\mid m\ \text{and}\ c\mid n,
\ \text{then}\ c\mid d.
\]
\[
m\mathbb Z\subseteq d\mathbb Z,\qquad
m\in m\mathbb Z\ \Longrightarrow\ (\exists)\,q\in\mathbb Z
\ \text{such that}\ m=dq,\quad d\mid m.
\]
\noindent Similarly, $d\mid n$ $\Longrightarrow d=(m,n)$.

\bigskip
\noindent $m\mathbb Z\cap n\mathbb Z\leq\mathbb Z\ \Longrightarrow\
(\exists!)\,u\in\mathbb N$ such that
\[
m\mathbb Z\cap n\mathbb Z=u\mathbb Z.
\]
\noindent $u\in m\mathbb Z$ and $u\in n\mathbb Z$
$\Longrightarrow m\mid u$ and $n\mid u$.

\noindent If $m\mid b$ and $n\mid b$,
$\stackrel{?}{\Longrightarrow} b\mid u$.
% The handwritten final relation appears to be b|u, although the
% subsequent membership statement implies u|b. Check the original.
\[
\underbrace{b\in m\mathbb Z}_{m\mid b},\quad
\underbrace{b\in n\mathbb Z}_{n\mid b}
\quad\Longrightarrow\quad b\in u\mathbb Z.
\]

\bigskip
\noindent First Isomorphism Theorem $\Longrightarrow$
\[
\frac{m\mathbb Z+n\mathbb Z}{n\mathbb Z}
\cong \frac{m\mathbb Z}{m\mathbb Z\cap n\mathbb Z}.
\]
\[
\Longrightarrow\qquad
\underbrace{\frac{d\mathbb Z}{n\mathbb Z}}_{\cong\,\mathbb Z_{n_1}}
\cong
\underbrace{\frac{m\mathbb Z}{u\mathbb Z}}_{\cong\,\mathbb Z_{u_1}}.
\]
% The n_1 and u_1 subscripts are slightly faded in the manuscript.
\[
u=m u_1,\qquad n=d n_1.
\]
\[
u_1=n_1\quad\Longrightarrow\quad
\frac{u}{m}=\frac{n}{d}\quad\Longrightarrow\quad du=mn.
\]
.
% No preamble or forced page break.
% Faithful reconstruction; uncertain readings are recorded in comments.

\subsection*{Lecture 11 - Page 6}

% Continuation of the isomorphism theorem proof from page 5.
\noindent $\underbrace{\bar a}_{aH}\in\ker f^*
\iff aH=aN\iff a\in H$.
% In the first line, aH=aN and a\in H are the apparent readings,
% but their relationship to the preceding homomorphism is unclear.
\[
\ker f^*=\{x\mid (\ldots)\}=\frac{H}{N}.
\]
% The expression after the vertical bar is partly cancelled/illegible.
% H/N appears to be written in the original. It has not been silently
% corrected to an alternative mathematical expression.

\bigskip
\noindent\underline{Applications:} \quad $(\mathbb Z,+,0)$.

\medskip
\noindent $H\leq\mathbb Z\ \Longrightarrow\ (\exists!)\,u\in\mathbb N^*$
 such that $H=u\mathbb Z$.
% The manuscript appears to specify N^*. Taken literally this excludes
% the trivial subgroup H={0}, for which the generator is 0.

\bigskip
\noindent Let $m,n\in\mathbb N^*$. I show that
\[
mn=[m,n]\cdot(m,n).
\]
% Brackets denote least common multiple; parentheses greatest common divisor.
\[
m\mathbb Z+n\mathbb Z\leq\mathbb Z
\ \Longrightarrow\ (\exists!)\,d\in\mathbb N\ \text{such that}\
m\mathbb Z+n\mathbb Z=d\mathbb Z.
\]
\noindent I show that $d=(m,n)$.

\medskip
\noindent There exist $x,y\in\mathbb Z$ such that
\[
mx+ny=d
\quad\Longrightarrow\quad
\text{if } c\mid m\ \text{and}\ c\mid n,
\ \text{then}\ c\mid d.
\]
\[
m\mathbb Z\subseteq d\mathbb Z,\qquad
m\in m\mathbb Z\ \Longrightarrow\ (\exists)\,q\in\mathbb Z
\ \text{such that}\ m=dq,\quad d\mid m.
\]
\noindent Similarly, $d\mid n$ $\Longrightarrow d=(m,n)$.

\bigskip
\noindent $m\mathbb Z\cap n\mathbb Z\leq\mathbb Z\ \Longrightarrow\
(\exists!)\,u\in\mathbb N$ such that
\[
m\mathbb Z\cap n\mathbb Z=u\mathbb Z.
\]
\noindent $u\in m\mathbb Z$ and $u\in n\mathbb Z$
$\Longrightarrow m\mid u$ and $n\mid u$.

\noindent If $m\mid b$ and $n\mid b$,
$\stackrel{?}{\Longrightarrow} b\mid u$.
% The handwritten final relation appears to be b|u, although the
% subsequent membership statement implies u|b. Check the original.
\[
\underbrace{b\in m\mathbb Z}_{m\mid b},\quad
\underbrace{b\in n\mathbb Z}_{n\mid b}
\quad\Longrightarrow\quad b\in u\mathbb Z.
\]

\bigskip
\noindent First Isomorphism Theorem $\Longrightarrow$
\[
\frac{m\mathbb Z+n\mathbb Z}{n\mathbb Z}
\cong \frac{m\mathbb Z}{m\mathbb Z\cap n\mathbb Z}.
\]
\[
\Longrightarrow\qquad
\underbrace{\frac{d\mathbb Z}{n\mathbb Z}}_{\cong\,\mathbb Z_{n_1}}
\cong
\underbrace{\frac{m\mathbb Z}{u\mathbb Z}}_{\cong\,\mathbb Z_{u_1}}.
\]
% The n_1 and u_1 subscripts are slightly faded in the manuscript.
\[
u=m u_1,\qquad n=d n_1.
\]
\[
u_1=n_1\quad\Longrightarrow\quad
\frac{u}{m}=\frac{n}{d}\quad\Longrightarrow\quad du=mn.
\]
.
% No preamble or forced page break.
% Faithful reconstruction; uncertain readings are recorded in comments.

\subsection*{Lecture 11 - Page 6}

% Continuation of the isomorphism theorem proof from page 5.
\noindent $\underbrace{\bar a}_{aH}\in\ker f^*
\iff aH=aN\iff a\in H$.
% In the first line, aH=aN and a\in H are the apparent readings,
% but their relationship to the preceding homomorphism is unclear.
\[
\ker f^*=\{x\mid (\ldots)\}=\frac{H}{N}.
\]
% The expression after the vertical bar is partly cancelled/illegible.
% H/N appears to be written in the original. It has not been silently
% corrected to an alternative mathematical expression.

\bigskip
\noindent\underline{Applications:} \quad $(\mathbb Z,+,0)$.

\medskip
\noindent $H\leq\mathbb Z\ \Longrightarrow\ (\exists!)\,u\in\mathbb N^*$
 such that $H=u\mathbb Z$.
% The manuscript appears to specify N^*. Taken literally this excludes
% the trivial subgroup H={0}, for which the generator is 0.

\bigskip
\noindent Let $m,n\in\mathbb N^*$. I show that
\[
mn=[m,n]\cdot(m,n).
\]
% Brackets denote least common multiple; parentheses greatest common divisor.
\[
m\mathbb Z+n\mathbb Z\leq\mathbb Z
\ \Longrightarrow\ (\exists!)\,d\in\mathbb N\ \text{such that}\
m\mathbb Z+n\mathbb Z=d\mathbb Z.
\]
\noindent I show that $d=(m,n)$.

\medskip
\noindent There exist $x,y\in\mathbb Z$ such that
\[
mx+ny=d
\quad\Longrightarrow\quad
\text{if } c\mid m\ \text{and}\ c\mid n,
\ \text{then}\ c\mid d.
\]
\[
m\mathbb Z\subseteq d\mathbb Z,\qquad
m\in m\mathbb Z\ \Longrightarrow\ (\exists)\,q\in\mathbb Z
\ \text{such that}\ m=dq,\quad d\mid m.
\]
\noindent Similarly, $d\mid n$ $\Longrightarrow d=(m,n)$.

\bigskip
\noindent $m\mathbb Z\cap n\mathbb Z\leq\mathbb Z\ \Longrightarrow\
(\exists!)\,u\in\mathbb N$ such that
\[
m\mathbb Z\cap n\mathbb Z=u\mathbb Z.
\]
\noindent $u\in m\mathbb Z$ and $u\in n\mathbb Z$
$\Longrightarrow m\mid u$ and $n\mid u$.

\noindent If $m\mid b$ and $n\mid b$,
$\stackrel{?}{\Longrightarrow} b\mid u$.
% The handwritten final relation appears to be b|u, although the
% subsequent membership statement implies u|b. Check the original.
\[
\underbrace{b\in m\mathbb Z}_{m\mid b},\quad
\underbrace{b\in n\mathbb Z}_{n\mid b}
\quad\Longrightarrow\quad b\in u\mathbb Z.
\]

\bigskip
\noindent First Isomorphism Theorem $\Longrightarrow$
\[
\frac{m\mathbb Z+n\mathbb Z}{n\mathbb Z}
\cong \frac{m\mathbb Z}{m\mathbb Z\cap n\mathbb Z}.
\]
\[
\Longrightarrow\qquad
\underbrace{\frac{d\mathbb Z}{n\mathbb Z}}_{\cong\,\mathbb Z_{n_1}}
\cong
\underbrace{\frac{m\mathbb Z}{u\mathbb Z}}_{\cong\,\mathbb Z_{u_1}}.
\]
% The n_1 and u_1 subscripts are slightly faded in the manuscript.
\[
u=m u_1,\qquad n=d n_1.
\]
\[
u_1=n_1\quad\Longrightarrow\quad
\frac{u}{m}=\frac{n}{d}\quad\Longrightarrow\quad du=mn.
\]
